\documentclass[aspectratio=169,11pt]{beamer}
% Shared CMPSC 480 Beamer preamble.
\usepackage[T1]{fontenc}
\usepackage[utf8]{inputenc}
\usepackage{lmodern}
\usepackage{amsmath,amssymb,mathtools}
\usepackage{booktabs,array,tabularx}
\usepackage{xcolor}
\usepackage{listings}
\usepackage{tikz}
\usetikzlibrary{arrows.meta,positioning,shapes.geometric}

% Ignore only negligible TeX box diagnostics that are below visual tolerance.
% Larger layout defects still appear in the build log and in rendered QA.
\vfuzz=3pt
\hbadness=2000

\definecolor{courseink}{HTML}{111111}
\definecolor{coursegray}{HTML}{EDEDED}
\definecolor{courserule}{HTML}{B8BCC4}
\definecolor{courseblue}{HTML}{3D8DFF}
\definecolor{courselightblue}{HTML}{D0EDFA}
\definecolor{coursemuted}{HTML}{5C6470}
\definecolor{coursegreen}{HTML}{0B7A53}
\definecolor{coursered}{HTML}{B3261E}

\usetheme{default}
\usefonttheme{professionalfonts}
\setbeamertemplate{navigation symbols}{}
\setbeamersize{text margin left=0.65cm,text margin right=0.65cm}

\setbeamercolor{normal text}{fg=courseink,bg=white}
\setbeamercolor{frametitle}{fg=courseink,bg=white}
\setbeamercolor{title}{fg=courseink,bg=white}
\setbeamercolor{subtitle}{fg=coursemuted,bg=white}
\setbeamercolor{structure}{fg=courseblue}
\setbeamercolor{alerted text}{fg=coursered}
\setbeamercolor{block title}{fg=courseink,bg=courselightblue}
\setbeamercolor{block body}{fg=courseink,bg=coursegray}
\setbeamercolor{example text}{fg=coursegreen}

\setbeamerfont{title}{size=\fontsize{25}{29}\selectfont,series=\bfseries}
\setbeamerfont{subtitle}{size=\fontsize{13}{16}\selectfont}
\setbeamerfont{frametitle}{size=\fontsize{19}{22}\selectfont,series=\bfseries}
\setbeamerfont{normal text}{size=\fontsize{11}{14}\selectfont}
\setbeamerfont{footline}{size=\fontsize{7.5}{9}\selectfont}

\setbeamertemplate{frametitle}{%
  \nointerlineskip
  % Use content-driven height so a long assessment title can wrap without
  % being clipped by a fixed-height title bar.
  \begin{beamercolorbox}[wd=\paperwidth,sep=0.17cm,leftskip=0.48cm,rightskip=0.48cm]{frametitle}
    \insertframetitle
  \end{beamercolorbox}
  \vspace{-0.08cm}
  {\color{courserule}\rule{\textwidth}{0.35pt}}
}

\setbeamertemplate{footline}{%
  \leavevmode
  \begin{beamercolorbox}[wd=\paperwidth,ht=2.6ex,dp=1.1ex,leftskip=.65cm,rightskip=.65cm]{author in head/foot}
    \color{coursemuted}\insertshorttitle\hfill\insertframenumber/\inserttotalframenumber
  \end{beamercolorbox}
}

\setbeamertemplate{itemize item}{\color{courseblue}\small$\blacktriangleright$}
\setbeamertemplate{itemize subitem}{\color{courseblue}\scriptsize$\bullet$}
\setbeamertemplate{enumerate item}{\color{courseblue}\insertenumlabel.}

\lstdefinestyle{coursepython}{
  language=Python,
  basicstyle=\ttfamily\fontsize{8.5}{10}\selectfont,
  keywordstyle=\color{courseblue}\bfseries,
  commentstyle=\color{coursemuted}\itshape,
  stringstyle=\color{coursegreen},
  showstringspaces=false,
  columns=fullflexible,
  keepspaces=true,
  frame=single,
  rulecolor=\color{courserule},
  backgroundcolor=\color{coursegray},
  breaklines=true,
  tabsize=4,
  xleftmargin=0.15cm,
  xrightmargin=0.15cm,
  framexleftmargin=0.15cm,
  framexrightmargin=0.15cm
}
\lstset{style=coursepython}

\newcommand{\points}[1]{\hfill{\color{courseblue}\textbf{[#1 points]}}}
\newcommand{\deliverable}[1]{\textbf{\color{courseblue}#1}}
\newcommand{\code}[1]{\texttt{#1}}
\newcommand{\keyidea}[1]{%
  \begin{beamercolorbox}[rounded=false,sep=0.55em,wd=\linewidth]{block body}
    \textbf{Key idea.} #1
  \end{beamercolorbox}
}
\newcommand{\answerlabel}{\textbf{\color{coursegreen}Answer.}\ }
\newcommand{\gradinglabel}{\textbf{\color{courseblue}Grading guidance.}\ }

\newenvironment{questionblock}[1]
  {\begin{block}{#1}}
  {\end{block}}

\newenvironment{solutionblock}[1]
  {\begin{block}{#1}}
  {\end{block}}

\AtBeginSection[]{
  \begin{frame}
    \vfill
    {\usebeamerfont{title}\color{courseink}\insertsectionhead\par}
    \vspace{0.4cm}
    {\color{courseblue}\rule{0.32\paperwidth}{3pt}}
    \vfill
  \end{frame}
}


% CMPSC480_TOTAL_POINTS: 20
% CMPSC480_ITEM_IDS: 1,2,3,4
\title[Quiz 1 Solutions]{Quiz 1: Agents and Search}
\subtitle{Weeks 1-2}
\author{CMPSC 480 - Instructor Solution Deck}
\date{}

\begin{document}


\begin{frame}
  \titlepage
\vspace{-0.35cm}
\begin{center}
\color{coursered}\bfseries Instructor material - do not distribute with the question deck
\end{center}
\end{frame}

\begin{frame}{Solution overview}
\begin{columns}[T,onlytextwidth]
\begin{column}{0.62\textwidth}
\Large This instructor deck provides a complete matched solution and grading guidance.

\vspace{0.35cm}
\normalsize
Coverage: Weeks 1-2
\end{column}
\begin{column}{0.33\textwidth}
\begin{block}{Points}20 total\end{block}
\begin{block}{Expected time}30 minutes\end{block}
\begin{block}{Submission}Individual PDF\end{block}
\end{column}
\end{columns}
\end{frame}

\begin{frame}{Instructor verification checklist}
\small
\begin{itemize}
  \item Work independently.
  \item Show intermediate frontier or arithmetic steps.
  \item State every assumption used for a guarantee.
  \item No electronic communication.
\end{itemize}
\end{frame}

\begin{frame}{Solution 1.a - Agent formulation (2 points)}
\small
\textbf{Question focus.} Give one element for each PEAS component.

\vspace{0.25cm}
\answerlabel Performance: timely collision-free deliveries; environment: aisles, shelves, people, and stations; actuators: wheel and lift commands; sensors: lidar, odometry, and task messages.
\end{frame}

\begin{frame}{Solution 1.b - Agent formulation (2 points)}
\small
\textbf{Question focus.} Explain why rational does not mean omniscient.

\vspace{0.25cm}
\answerlabel A rational agent maximizes expected performance from its percept history and knowledge. It can still fail because observations and outcomes are uncertain; omniscience would require knowing actual future consequences.
\end{frame}

\begin{frame}{Solution 2.a - BFS and DFS trace (2 points)}
\small
\textbf{Question focus.} Trace BFS until a goal is removed.

\vspace{0.25cm}
\answerlabel FIFO trace: [S], then [A,B], then [B,G], then [G,C], then remove G. The returned path is S-A-G.
\end{frame}

\begin{frame}{Solution 2.b - BFS and DFS trace (2 points)}
\small
\textbf{Question focus.} Trace DFS when successors are pushed in alphabetical order.

\vspace{0.25cm}
\answerlabel After pushing A then B, LIFO removes B, then C, then G; it can return S-B-C-G. A different documented push convention earns credit if traced consistently.
\end{frame}

\begin{frame}{Solution 3.a - Cost-sensitive search (2 points)}
\small
\textbf{Question focus.} Which route can BFS return, and why is that not a minimum-cost guarantee?

\vspace{0.25cm}
\answerlabel BFS returns the depth-two route S-A-G with cost 6, because it minimizes depth under FIFO order rather than nonuniform path cost.
\end{frame}

\begin{frame}{Solution 3.b - Cost-sensitive search (2 points)}
\small
\textbf{Question focus.} Trace UCS priorities until the goal is removed.

\vspace{0.25cm}
\answerlabel UCS removes S(0), B(1), C(2), G(3), and returns S-B-C-G with cost 3.
\end{frame}

\begin{frame}{Solution 3.c - Cost-sensitive search (2 points)}
\small
\textbf{Question focus.} State conditions under which UCS is complete and optimal.

\vspace{0.25cm}
\answerlabel With finite branching and step costs bounded below by a positive epsilon, UCS is complete and removes an optimal goal first.
\end{frame}

\begin{frame}{Solution 4.a - A-star heuristic (2 points)}
\small
\textbf{Question focus.} Define f, g, and h in A-star.

\vspace{0.25cm}
\answerlabel g is accumulated path cost, h estimates optimal remaining cost, and f=g+h orders the frontier.
\end{frame}

\begin{frame}{Solution 4.b - A-star heuristic (2 points)}
\small
\textbf{Question focus.} Prove Manhattan distance is consistent.

\vspace{0.25cm}
\answerlabel A legal step costs one and changes Manhattan distance by at most one, hence h(n) is at most 1+h(n-prime) for every edge.
\end{frame}

\begin{frame}{Solution 4.c - A-star heuristic (2 points)}
\small
\textbf{Question focus.} If two nodes have (g,h)=(4,3) and (5,1), which is removed first?

\vspace{0.25cm}
\answerlabel Their f values are 7 and 6, so the second node is removed first, subject only to stale-entry handling and ties.
\end{frame}

\begin{frame}{Edge-case reasoning and expected behavior}
\small
\begin{itemize}
  \item Return the empty path.
  \item Use graph-search bookkeeping to terminate.
  \item Apply the stated deterministic secondary key.
\end{itemize}
\end{frame}

\begin{frame}{Grading rubric}
\scriptsize
\begin{tabularx}{\textwidth}{>{\bfseries}p{0.28\textwidth} p{0.08\textwidth} X}
\toprule
Category & Points & Full-credit evidence \\
\midrule
Agent formulation & 4 & PEAS and rationality are precise. \\
BFS and DFS trace & 4 & FIFO/LIFO order and trace are correct. \\
Cost-sensitive search & 6 & Depth/cost distinction and UCS update are correct. \\
A-star heuristic & 6 & Admissibility, consistency, and f values are correct. \\
\bottomrule
\end{tabularx}
\end{frame}

\begin{frame}{Reflection exemplar}
\small
\answerlabel A suitable answer identifies finite branching, uniform costs for BFS optimality, or a positive lower bound on step cost for UCS, with the relevant algorithm named.

\vspace{0.25cm}
\gradinglabel Award credit for a specific claim supported by technical evidence from the submitted work.
\end{frame}

\end{document}
